\section{Título del Informe}

Informe de Evaluación de Usabilidad y Satisfacción del Usuario - [Nombre de la Aplicación]

\section{Introducción}

En este informe se presenta la evaluación de usabilidad y satisfacción del usuario de la aplicación [Nombre de la Aplicación]. El objetivo principal es analizar la facilidad de uso y el grado de satisfacción de los usuarios mediante la aplicación de pruebas con usuarios reales, encuestas y entrevistas. La metodología empleada incluye la observación directa de sesiones de uso, la recopilación de feedback y el análisis de métricas de satisfacción.

\section{Plan de Pruebas de Usabilidad}

\subsection{Perfiles de usuarios participantes}
Se definieron perfiles de usuarios representativos del público objetivo de la aplicación, considerando diferentes niveles de experiencia y necesidades.

\subsection{Escenarios y tareas de prueba}
Se diseñaron escenarios y tareas específicas para evaluar la interacción de los usuarios con las funcionalidades principales de la aplicación.

\subsection{Cuestionarios utilizados}
Se utilizó el cuestionario SUS (System Usability Scale) para medir la percepción de usabilidad y satisfacción de los usuarios.

\section{Hallazgos de Usabilidad}

Durante las sesiones de prueba se identificaron los siguientes problemas de usabilidad:
\begin{itemize}
    \item [Ejemplo] Dificultad para encontrar la opción de registro.
    \item [Ejemplo] Interfaz poco intuitiva en el proceso de navegación.
\end{itemize}

Las métricas de satisfacción obtenidas a través del cuestionario SUS arrojaron una puntuación promedio de [XX/100].

\section{Recomendaciones de Mejora}

Para cada problema de usabilidad identificado, se proponen las siguientes soluciones concretas y priorizadas:
\begin{itemize}
    \item [Ejemplo] Mejorar la visibilidad del botón de registro.
    \item [Ejemplo] Simplificar la estructura de menús para facilitar la navegación.
\end{itemize}

Se espera que la implementación de estas mejoras tenga un impacto positivo en la experiencia del usuario, incrementando la satisfacción y eficiencia en el uso de la aplicación.

\section{Conclusiones}

En resumen, la evaluación realizada evidencia el nivel de usabilidad y satisfacción del usuario en la aplicación [Nombre de la Aplicación]. Se recomienda implementar las mejoras propuestas para optimizar la experiencia del usuario.

\section*{Evidencias}
\begin{itemize}
    \item Clips de video grabados con OBS Studio
    \item Capturas de pantalla de la interfaz
    \item Gráficos de resultados SUS
\end{itemize}
